\pagestyle{empty}
\tight
\begin{tblr}{width=\textwidth,colspec={|Q[c,m,1.9cm]|X[l,m]|},hlines,vlines,hspan=minimal,row{1}={bg=headblue},rowsep=1pt,colsep=3pt}
\SetCell{c,m}\hfil\logo\hfil & \textbf{POLITEKNIK NEGERI MALANG}\par\textbf{JURUSAN TEKNOLOGI INFORMASI}\par\textbf{PROGRAM STUDI: D4 TEKNIK INFORMATIKA} \\
\SetCell[c=2]{c} \textbf{RENCANA PEMBELAJARAN SEMESTER (RPS)} & \\
\end{tblr}
\begin{longtblr}{width=\textwidth,colspec={|Q[l,m,1.9cm]|X[0.85,l,m]|X[1.45,l,m]|X[0.75,l,m]|X[1.15,l,m]|},hlines,vlines,hspan=minimal,rowsep=1pt,colsep=2pt,row{1,3}={bg=headgray,font=\bfseries}}
MATA KULIAH & KODE & BOBOT (SKS)/jam & SEMESTER & TGL. PENYUSUNAN \\
Proyek Inovasi & RTI256204 & 6 SKS/12 jam & 6 & 1 Juli 2025 \\
OTORISASI & Dosen Pengembang RPS & Koordinator RMK & \SetCell[c=2]{l} Ka PRODI &  \\
 & Vivi Nur Wijayaningrum, S.Kom., M.Kom. &  & \SetCell[c=2]{l}{\vspace{8mm}\par Dr. Ely Setyo Astuti, S.T., M.T.} &  \\
\SetCell[r=13]{l} Capaian Pembelajaran (CP) & \SetCell[c=4]{l,bg=headgray,font=\bfseries} Capaian Pembelajaran Lulusan Program Studi (CPL-Prodi) &  &  &  \\
 & CPL08 & \SetCell[c=3]{l} Mampu mengelola proyek teknologi informasi secara profesional melalui perencanaan, koordinasi, komunikasi, dan optimalisasi sumber daya. &  &  \\
 & CPL06 & \SetCell[c=3]{l} Mampu merancang, mengimplementasikan, menguji, melakukan deployment, memelihara, serta menjamin mutu perangkat lunak yang menjawab permasalahan dan memenuhi kebutuhan pengguna. &  &  \\
 & \SetCell[c=4]{l,bg=headgray,font=\bfseries} Capaian Pembelajaran mata kuliah (CPL-MK) &  &  &  \\
 & CPMK0801 & \SetCell[c=3]{l} Mahasiswa mampu merencanakan proyek inovasi TI berbasis konteks industri nyata. &  &  \\
 & CPMK0803 & \SetCell[c=3]{l} Mahasiswa mampu mengeksekusi, memonitor, dan mengendalikan proyek inovasi serta memimpin tim. &  &  \\
 & CPMK0804 & \SetCell[c=3]{l} Mahasiswa mampu mendemonstrasikan penyelesaian proyek inovasi TI secara utuh. &  &  \\
 & CPMK0607 & \SetCell[c=3]{l} Mahasiswa mampu menghasilkan solusi sistem informasi inovatif yang menjawab permasalahan nyata industri. &  &  \\
 & \SetCell[c=4]{l,bg=headgray,font=\bfseries} Kemampuan akhir tiap tahapan belajar (Sub-CPMK) &  &  &  \\
 & SCPMK0801-05301 & \SetCell[c=3]{l} Mahasiswa mampu menyusun proposal proyek inovasi TI dengan analisis masalah, solusi teknologi, WBS, dan rencana risiko. &  &  \\
 & SCPMK0803-05302 & \SetCell[c=3]{l} Mahasiswa mampu mengeksekusi proyek inovasi berbasis sprint, melaporkan kemajuan, dan mengelola perubahan tim. &  &  \\
 & SCPMK0804-05303 & \SetCell[c=3]{l} Mahasiswa mampu mempresentasikan hasil proyek inovasi dengan defense teknis dan bisnis yang komprehensif. &  &  \\
 & SCPMK0607-05304 & \SetCell[c=3]{l} Mahasiswa mampu mengintegrasikan teknologi mutakhir menjadi solusi inovatif yang scalable dan deployable. &  &  \\
\SetCell[r=2]{l} Penilaian & \SetCell[c=4]{l} *Nilai CPMK didapat dari agregasi nilai SUB-CPMK yang dibebankan &  &  &  \\
 & \SetCell[c=4]{l}{\tiny\begin{tblr}{width=\linewidth,colspec={|Q[c,m,0.1\linewidth]|Q[c,m,0.16\linewidth]|X[l,m]|X[l,m]|X[l,m]|X[l,m]|X[l,m]|X[l,m]|Q[c,m,0.08\linewidth]|},hlines,vlines,hspan=minimal,rowsep=1pt,colsep=0.7pt,row{1,2}={bg=headgray,font=\bfseries}}
Kode CPMK & Kode SCPMK & \SetCell[c=6]{c} Bobot per Bentuk Penilaian & & & & & & Total Bobot \\
 & & Proposal & Sprint 1-2 & UTS Milestone & Sprint 3-4 & Laporan & UAS Defense & \\
CPMK0801 & SCPMK0801-05301 & 15\%: proposal inovasi TI & 0\% & 0\% & 0\% & 0\% & 0\% & 15\% \\
CPMK0803 & SCPMK0803-05302 & 0\% & 10\%: sprint review & 15\%: milestone & 0\% & 0\% & 0\% & 25\% \\
CPMK0804 & SCPMK0804-05303 & 0\% & 0\% & 0\% & 10\%: sprint 3-4 & 10\%: laporan akhir & 10\%: defense & 30\% \\
CPMK0607 & SCPMK0607-05304 & 0\% & 10\%: implementasi & 0\% & 10\%: implementasi & 0\% & 10\%: demo inovasi & 30\% \\
\SetCell[c=8]{r}\textbf{Total Per Penilaian} & & & & & & & & \textbf{100\%} \\
\end{tblr}} &  &  &  \\
Diskripsi Singkat Mata Kuliah & \SetCell[c=4]{l} Proyek Inovasi adalah capstone jalur magang/industri yang mengintegrasikan kreativitas, rekayasa, dan manajemen proyek dalam mengembangkan solusi TI inovatif berbasis masalah nyata industri menggunakan metodologi agile. &  &  &  \\
Bahan Kajian / Materi Pembelajaran & \SetCell[c=4]{l}\begin{enumerate}\item Problem framing, opportunity analysis, dan technology landscape\item Proposal inovasi: WBS, jadwal, anggaran, dan manajemen risiko\item Metodologi agile/scrum: sprint planning, review, dan retrospective\item Implementasi MVP, UAT, deployment, dan dokumentasi proyek\end{enumerate} &  &  &  \\
\SetCell[r=2]{l} Pustaka & \SetCell[c=4]{l} \textbf{Utama:}\begin{enumerate}\item Ries, E. 2011. The Lean Startup. Crown Business.\item Blank, S. 2013. The Startup Owner's Manual. K\&S Ranch.\item Schwaber, K. 2020. The Scrum Guide. Scrum.org.\end{enumerate} &  &  &  \\
 & \SetCell[c=4]{l} \textbf{Pendukung:} Materi LMS, dokumentasi resmi perangkat lunak, dan jobsheet praktikum. &  &  &  \\
Media Pembelajaran & \SetCell[c=2]{l} \textbf{Software:} Project management tools (Jira/Trello), Git, IDE, LMS. &  & \SetCell[c=2]{l} \textbf{Hardware:} Komputer/laptop, proyektor. &  \\
Nama Dosen Pengampu & \SetCell[c=4]{l} Tim Pengajar Program Studi D4 TI &  &  &  \\
Matakuliah Syarat & \SetCell[c=4]{l} - &  &  &  \\
\end{longtblr}
\begin{longtblr}{width=\textwidth,colspec={|Q[c,m,0.06\textwidth]|X[1.35,l,m]|X[1.08,l,m]|X[1.16,l,m]|X[0.7,l,m]|X[1.24,l,m]|X[1.06,l,m]|X[0.98,l,m]|Q[c,m,0.05\textwidth]|},hlines,vlines,hspan=minimal,rowhead=2,row{1,2}={bg=headgreen,font=\bfseries},rowsep=1pt,colsep=1.5pt}
Mgg.\par Ke & Kemampuan Akhir Yang Direncanakan (Sub-CP-MK) & Bahan kajian (Materi Pembelajaran) & Bentuk dan Metode Pembelajaran & Estimasi Waktu & Pengalaman Belajar Mahasiswa & Kriteria \& Bentuk Penilaian & Indikator Penilaian & Bobot Penilaian (\%) \\
(1) & (2) & (3) & (4) & (5) & (6) & (7) & (8) & (9) \\
1 & SCPMK0801-05301 & Orientasi proyek inovasi dan pembentukan tim. & Modalitas: Blended Learning. Bentuk: Luring/Daring. Metode: Project Based Learning, agile sprint, presentasi, dan peer review. & 1 x 4 x 50' tatap muka; 1 x 4 x 50' proyek mandiri. & Mahasiswa mengerjakan aktivitas terarah untuk orientasi proyek inovasi dan menunjukkan evidence hasil belajar. & Bentuk penilaian: Orientasi Proyek Inovasi. Mengacu rubrik RTM. & Kelengkapan artefak; kualitas implementasi; kemajuan dan dokumentasi. & 3 \\
2 & SCPMK0801-05301 & Problem framing dan opportunity analysis. & Modalitas: Blended Learning. Bentuk: Luring/Daring. Metode: Project Based Learning, agile sprint, presentasi, dan peer review. & 1 x 4 x 50' tatap muka; 1 x 4 x 50' proyek mandiri. & Mahasiswa mengerjakan aktivitas terarah untuk problem framing dan opportunity analysis dan menunjukkan evidence hasil belajar. & Bentuk penilaian: Problem Framing dan Opportunity Analysis. Mengacu rubrik RTM. & Kelengkapan artefak; kualitas implementasi; kemajuan dan dokumentasi. & 3 \\
3 & SCPMK0801-05301 & Technology landscape: tren TI 2025. & Modalitas: Blended Learning. Bentuk: Luring/Daring. Metode: Project Based Learning, agile sprint, presentasi, dan peer review. & 1 x 4 x 50' tatap muka; 1 x 4 x 50' proyek mandiri. & Mahasiswa mengerjakan aktivitas terarah untuk technology landscape dan menunjukkan evidence hasil belajar. & Bentuk penilaian: Technology Landscape. Mengacu rubrik RTM. & Kelengkapan artefak; kualitas implementasi; kemajuan dan dokumentasi. & 3 \\
4 & SCPMK0801-05301 & Proposal: WBS, jadwal, anggaran, risiko. & Modalitas: Blended Learning. Bentuk: Luring/Daring. Metode: Project Based Learning, agile sprint, presentasi, dan peer review. & 1 x 4 x 50' tatap muka; 1 x 4 x 50' proyek mandiri. & Mahasiswa mengerjakan aktivitas terarah untuk penyusunan proposal dan menunjukkan evidence hasil belajar. & Bentuk penilaian: Proposal Proyek Inovasi. Mengacu rubrik RTM. & Kelengkapan artefak; kualitas implementasi; kemajuan dan dokumentasi. & 6 \\
5 & SCPMK0607-05304 & Sprint 1: Proof of concept dan MVP. & Modalitas: Blended Learning. Bentuk: Luring/Daring. Metode: Project Based Learning, agile sprint, presentasi, dan peer review. & 1 x 4 x 50' tatap muka; 1 x 4 x 50' proyek mandiri. & Mahasiswa mengerjakan aktivitas terarah untuk sprint 1 proof of concept dan menunjukkan evidence hasil belajar. & Bentuk penilaian: Sprint 1 PoC dan MVP. Mengacu rubrik RTM. & Kelengkapan artefak; kualitas implementasi; kemajuan dan dokumentasi. & 4 \\
6 & SCPMK0607-05304 & Sprint 1: validasi teknis dan iterasi. & Modalitas: Blended Learning. Bentuk: Luring/Daring. Metode: Project Based Learning, agile sprint, presentasi, dan peer review. & 1 x 4 x 50' tatap muka; 1 x 4 x 50' proyek mandiri. & Mahasiswa mengerjakan aktivitas terarah untuk sprint 1 validasi teknis dan menunjukkan evidence hasil belajar. & Bentuk penilaian: Sprint 1 Validasi Teknis. Mengacu rubrik RTM. & Kelengkapan artefak; kualitas implementasi; kemajuan dan dokumentasi. & 4 \\
7 & SCPMK0803-05302 & Sprint 1 review dan demo. & Modalitas: Blended Learning. Bentuk: Luring/Daring. Metode: Project Based Learning, agile sprint, presentasi, dan peer review. & 1 x 4 x 50' tatap muka; 1 x 4 x 50' proyek mandiri. & Mahasiswa mengerjakan aktivitas terarah untuk sprint 1 review dan demo dan menunjukkan evidence hasil belajar. & Bentuk penilaian: Sprint 1 Review dan Demo. Mengacu rubrik RTM. & Kelengkapan artefak; kualitas implementasi; kemajuan dan dokumentasi. & 5 \\
8 & SCPMK0803-05302 & UTS Milestone review dan evaluasi. & Modalitas: Blended Learning. Bentuk: Luring/Daring. Metode: Project Based Learning, agile sprint, presentasi, dan peer review. & 1 x 4 x 50' tatap muka; 1 x 4 x 50' proyek mandiri. & Mahasiswa mengerjakan aktivitas terarah untuk UTS milestone review dan menunjukkan evidence hasil belajar. & Bentuk penilaian: UTS Milestone Review. Mengacu rubrik RTM. & Kelengkapan artefak; kualitas implementasi; kemajuan dan dokumentasi. & 15 \\
9 & SCPMK0607-05304 & Sprint 2: pengembangan fitur utama. & Modalitas: Blended Learning. Bentuk: Luring/Daring. Metode: Project Based Learning, agile sprint, presentasi, dan peer review. & 1 x 4 x 50' tatap muka; 1 x 4 x 50' proyek mandiri. & Mahasiswa mengerjakan aktivitas terarah untuk sprint 2 pengembangan fitur utama dan menunjukkan evidence hasil belajar. & Bentuk penilaian: Sprint 2 Pengembangan Fitur. Mengacu rubrik RTM. & Kelengkapan artefak; kualitas implementasi; kemajuan dan dokumentasi. & 4 \\
10 & SCPMK0607-05304 & Sprint 2: integration testing dan optimasi. & Modalitas: Blended Learning. Bentuk: Luring/Daring. Metode: Project Based Learning, agile sprint, presentasi, dan peer review. & 1 x 4 x 50' tatap muka; 1 x 4 x 50' proyek mandiri. & Mahasiswa mengerjakan aktivitas terarah untuk sprint 2 integration testing dan menunjukkan evidence hasil belajar. & Bentuk penilaian: Sprint 2 Integration Testing. Mengacu rubrik RTM. & Kelengkapan artefak; kualitas implementasi; kemajuan dan dokumentasi. & 5 \\
11 & SCPMK0803-05302 & Sprint 2 review dan Sprint 3 planning. & Modalitas: Blended Learning. Bentuk: Luring/Daring. Metode: Project Based Learning, agile sprint, presentasi, dan peer review. & 1 x 4 x 50' tatap muka; 1 x 4 x 50' proyek mandiri. & Mahasiswa mengerjakan aktivitas terarah untuk sprint 2 review dan sprint 3 planning dan menunjukkan evidence hasil belajar. & Bentuk penilaian: Sprint 2 Review dan Sprint 3 Planning. Mengacu rubrik RTM. & Kelengkapan artefak; kualitas implementasi; kemajuan dan dokumentasi. & 4 \\
12 & SCPMK0607-05304 & Sprint 3: skalabilitas dan deployment preparation. & Modalitas: Blended Learning. Bentuk: Luring/Daring. Metode: Project Based Learning, agile sprint, presentasi, dan peer review. & 1 x 4 x 50' tatap muka; 1 x 4 x 50' proyek mandiri. & Mahasiswa mengerjakan aktivitas terarah untuk sprint 3 skalabilitas dan deployment dan menunjukkan evidence hasil belajar. & Bentuk penilaian: Sprint 3 Skalabilitas dan Deployment. Mengacu rubrik RTM. & Kelengkapan artefak; kualitas implementasi; kemajuan dan dokumentasi. & 5 \\
13 & SCPMK0607-05304 & Sprint 3: UAT. & Modalitas: Blended Learning. Bentuk: Luring/Daring. Metode: Project Based Learning, agile sprint, presentasi, dan peer review. & 1 x 4 x 50' tatap muka; 1 x 4 x 50' proyek mandiri. & Mahasiswa mengerjakan aktivitas terarah untuk sprint 3 UAT dan menunjukkan evidence hasil belajar. & Bentuk penilaian: Sprint 3 UAT. Mengacu rubrik RTM. & Kelengkapan artefak; kualitas implementasi; kemajuan dan dokumentasi. & 5 \\
14 & SCPMK0804-05303 & Sprint 4: product polish dan dokumentasi. & Modalitas: Blended Learning. Bentuk: Luring/Daring. Metode: Project Based Learning, agile sprint, presentasi, dan peer review. & 1 x 4 x 50' tatap muka; 1 x 4 x 50' proyek mandiri. & Mahasiswa mengerjakan aktivitas terarah untuk sprint 4 product polish dan dokumentasi dan menunjukkan evidence hasil belajar. & Bentuk penilaian: Sprint 4 Product Polish. Mengacu rubrik RTM. & Kelengkapan artefak; kualitas implementasi; kemajuan dan dokumentasi. & 5 \\
15 & SCPMK0804-05303 & Penyusunan laporan akhir. & Modalitas: Blended Learning. Bentuk: Luring/Daring. Metode: Project Based Learning, agile sprint, presentasi, dan peer review. & 1 x 4 x 50' tatap muka; 1 x 4 x 50' proyek mandiri. & Mahasiswa mengerjakan aktivitas terarah untuk penyusunan laporan akhir dan menunjukkan evidence hasil belajar. & Bentuk penilaian: Laporan Akhir Proyek. Mengacu rubrik RTM. & Kelengkapan artefak; kualitas implementasi; kemajuan dan dokumentasi. & 5 \\
16 & SCPMK0804-05303 & Latihan presentasi dan pitch. & Modalitas: Blended Learning. Bentuk: Luring/Daring. Metode: Project Based Learning, agile sprint, presentasi, dan peer review. & 1 x 4 x 50' tatap muka; 1 x 4 x 50' proyek mandiri. & Mahasiswa mengerjakan aktivitas terarah untuk latihan presentasi dan pitch dan menunjukkan evidence hasil belajar. & Bentuk penilaian: Latihan Presentasi dan Pitch. Mengacu rubrik RTM. & Kelengkapan artefak; kualitas implementasi; kemajuan dan dokumentasi. & 4 \\
17 & SCPMK0804-05303, SCPMK0607-05304 & UAS Defense proyek inovasi dan demo. & Modalitas: Blended Learning. Bentuk: Luring/Daring. Metode: Project Based Learning, agile sprint, presentasi, dan peer review. & 1 x 4 x 50' tatap muka; 1 x 4 x 50' proyek mandiri. & Mahasiswa mengerjakan aktivitas terarah untuk UAS defense proyek inovasi dan menunjukkan evidence hasil belajar. & Bentuk penilaian: UAS Defense Proyek Inovasi. Mengacu rubrik RTM. & Kelengkapan artefak; kualitas implementasi; kemajuan dan dokumentasi. & 20 \\
\end{longtblr}
